\subsection{Separating source optimisation error from transfer error}
\label{sec:v24-selection-gap}
The V23 continuous-path enclosure constrains the maximum
\emph{source-validation} macro-AUC; it does not constrain the
performance of an unseen target participant. This distinction can
be made precise using a simple deterministic inequality.

Let $S(\lambda)$ be the participant-macro AUC of frozen
held-source validation scores, and $T(\lambda)$ the AUC on
a fixed unseen target participant, on a common feasible candidate
set $\Lambda$. Suppose a source-only selector chooses
$\widehat\lambda$ with
\[
 \max_{\lambda\in\Lambda} S(\lambda)
   -S(\widehat\lambda)\le\varepsilon_S.
\]
Define the unobservable-at-selection transfer discrepancy
\[
 \Delta_{S,T}(\Lambda)=
 \sup_{\lambda\in\Lambda}|S(\lambda)-T(\lambda)|.
\]
Then a direct add-and-subtract argument gives
\begin{equation}
 \underbrace{\max_{\lambda\in\Lambda}T(\lambda)
   -T(\widehat\lambda)}_{\text{target-oracle regret}}
 \le \underbrace{\varepsilon_S}_{\text{source selection gap}}
 +\underbrace{2\Delta_{S,T}(\Lambda)}_{\text{source--target discrepancy}}.
 \label{eq:v24-transfer-gap}
\end{equation}
This elementary inequality is not a newly discovered statistical
generalisation theorem. Its contribution to the study is to make
explicit \emph{why better source optimisation alone may be
irrelevant to transfer}: closing the source optimisation gap does
not control the second term. A source-only certificate for
$\varepsilon_S$ therefore must not be reported as a guarantee
on Eq.~\eqref{eq:v24-transfer-gap}; target labels would be required
to compute $\Delta_{S,T}$ post hoc. Non-exchangeable and label-reversed
target populations can have the same observed source data but
conflicting best shrinkage decisions. No source-only selector can
promise universally improved target AUC without additional
assumptions.

\paragraph{Matched source search repair.}
The original V23 exploratory continuous optimiser started with two
interval endpoints and five bisections and therefore sometimes
failed to evaluate the original 27-point grid's best candidate
within the same interval. The V24 repair first evaluates the
\emph{predeclared source-only} feasible finite-grid candidates,
partitions the continuous interval at them and searches
subintervals with pair-credit upper bounds. This eliminates
a logically avoidable finite-grid regression, at the cost of
additional candidate evaluation; that cost is reported.

In a genuine BNCI2014\_002 held-source experiment (first
outer pool, 13 validation source people), the original grid
and the repaired continuous search both attained source
macro-AUC $0.79397$ with selected $\lambda=0.8$.
The computed conditional source upper bound decreased to
$0.81019$, so the remaining source-only optimum gap
was at most $0.01623$ within the fixed interval
$[0.05,0.95]$ under the numerical enclosure assumptions.
This met the predeclared 0.02 \emph{source} AUC tolerance,
but yielded \emph{zero observed source AUC improvement}
relative to the already competitive grid. The earlier
endpoints-only five-split search had source AUC $0.79296$
and an inconclusive $0.09477$ source optimality gap.
This is a repaired optimisation audit, not a new target
transfer gain or clinical safety certificate.

\paragraph{External decision gate.}
A new public MI acquisition (Weibo2014, 10 participants,
three fixed sensorimotor electrodes and binary left/right
imagery) is reserved by a repository-committed, pre-outcome
source-only first-run protocol. Four source-only comparators
and a fifth fixed-grid control are chosen before target
scores are evaluated. Successful execution and actual
participant-level target outcomes are required before any
algorithm-performance statement can appear in this
manuscript; an in-repository freeze is not equivalent
to an independent registry preregistration. If the anchored
search matches the 27-point baseline without target benefit,
we will retain the negative result rather than claim the
continuous optimiser is a superior transfer algorithm.
